% Part: applied-modal-logics
% Chapter: epistemic-logic
% Section: public-announcement-logic-semantics

\documentclass[../../../include/open-logic-section]{subfiles}

\begin{document}

\olfileid{aml}{el}{psm}

\olsection{જાહેર જાહેરાતના તર્કનું અર્થવિજ્ઞાન}

જાહેર જાહેરાતના તર્ક માટેનાં સંબંધાત્મક નિદર્શો જ્ઞાનસંબંધી
તર્કનાં નિદર્શો જેવાં જ છે. પરંતુ જાહેર-જાહેરાત સંક્રિયકનું
અર્થવિજ્ઞાન નવું છે.

\begin{defn}\ollabel{defn:mmodels} નિદર્શ~$\mModel M = \tuple{W, R, V}$માં
  \emph{!!a{formula}~$!A$નું~$w$માં સત્ય હોવું}, સંજ્ઞામાં
  $\mSat{M}{!A}[w]$, નીચે મુજબ આગમનાત્મક રીતે વ્યાખ્યાયિત થાય છે:
  \begin{enumerate}
  \tagitem{prvFalse}{\indcase{!A}{\lfalse}{$\mSat{M}{\lfalse}[w]$ ક્યારેય નહીં}.}{}
  \tagitem{prvTrue}{\indcase{!A}{\ltrue}{$\mSat{M}{\ltrue}[w]$ હંમેશા}.}{}
  \item $\mSat{M}{p}[w]$ ત્યારે અને માત્ર ત્યારે જ જ્યારે $w \in V(p)$
  \tagitem{prvNot}{\indcase{!A}{\lnot !B}{$\mSat{M}{\indfrm}[w]$ ત્યારે અને માત્ર ત્યારે જ જ્યારે
    $\mSat/{M}{!B}[w]$}.}{}
  \tagitem{prvAnd}{\indcase{!A}{(!B \land !C)}{$\mSat{M}{\indfrm}[w]$ ત્યારે અને માત્ર ત્યારે જ જ્યારે
    $\mSat{M}{!B}[w]$ અને $\mSat{M}{!C}[w]$}.}{}
  \tagitem{prvOr}{\indcase{!A}{(!B \lor !C)}{$\mSat{M}{\indfrm}[w]$ ત્યારે અને માત્ર ત્યારે જ જ્યારે
    $\mSat{M}{!B}[w]$ અથવા $\mSat{M}{!C}[w]$} (અથવા બંને).}{}
  \tagitem{prvIf}{\indcase{!A}{(!B \lif !C)}{$\mSat{M}{\indfrm}[w]$ ત્યારે અને માત્ર ત્યારે જ જ્યારે
    $\mSat/{M}{!B}[w]$ અથવા $\mSat{M}{!C}[w]$}.}{}
  \tagitem{prvIff}{\indcase{!A}{(!B \liff !C)}{$\mSat{M}{\indfrm}[w]$ ત્યારે અને માત્ર ત્યારે જ જ્યારે
    $\mSat{M}{!B}[w]$ અને $\mSat{M}{!C}[w]$ બંને હોય, અથવા
    $\mSat{M}{!B}[w]$ તથા $\mSat{M}{!C}[w]$ બેમાંથી એકેય ન હોય}.}{}
  \item\ollabel{defn:sub:mmodels-box}
    \indcase{!A}{\Knows_a !B}{$\mSat{M}{\indfrm}[w]$ ત્યારે અને માત્ર ત્યારે જ જ્યારે
    $R_a ww'$ ધરાવતા દરેક $w' \in W$ માટે $\mSat{M}{!B}[w']$}
  \item\ollabel{defn:sub:mmodels-pal}
    \indcase{!A}{[!B] !C}{$\mSat{M}{\indfrm}[w]$ ત્યારે અને માત્ર ત્યારે જ જ્યારે
    $\mSat{M}{!B}[w]$ હોય તો $\mSat{M \mid !B}{!C}[w]$}

  
  અહીં $\mModel M \mid !B = \tuple{W', R', V'}$ નીચે મુજબ
  વ્યાખ્યાયિત થાય છે:
  \begin{enumerate}
    \item $W' = \Setabs{ u \in W}{\mSat{M}{!B}[u]}$. એટલે
      $\mModel M \mid !B$નાં જગતો, $\mModel M$માં $!B$ સત્ય હોય
      તે જ જગતો છે.
    \item $R'_a = R_a \cap (W' \times W')$. દરેક કર્તાનો સુલભતા
      સંબંધ ફક્ત~$W'$માં બચેલાં જગતો સુધી મર્યાદિત થાય છે.
    \item $V'(p) = \Setabs{ u \in W'}{u \in V(p) }$. એવી જ રીતે,
      જગતોમાં વિધાનાત્મક ચલની સત્ય-નિમણૂક એ જ રહે છે: માહિતીલક્ષી
      ઘટનાઓ વિધાનાત્મક ચલોનાં સત્યમૂલ્યો બદલતી નથી.
  \end{enumerate}
  
  \end{enumerate} 
\end{defn}

જાહેર જાહેરાતના તર્કની વિશિષ્ટતા એ છે કે નિદર્શ~$\mModel M$માં
!!a{formula}નું સત્ય હોવું ક્યારેક $\mModel M$ સિવાયના બીજા
નિદર્શને જોઈને જ નક્કી કરી શકાય છે.

આપણું અર્થવિજ્ઞાન જાહેરાતના સંક્રિયકને $\Box$~સંક્રિયક ગણે છે.
તેથી કોઈ જગતમાં !!a{formula}~$!A$ સત્યપૂર્વક જાહેર કરી ન શકાય,
તો ત્યાં $[!A]!B$ રિક્તપણે સત્ય હશે; જેમ અંતિમ જગતોમાં બધાં
$\Box$ !!{formula}s સત્ય હોય છે.
\sourcecorrection{OLMD-081}{મૂળમાં અહીં $[!A]B$ છે, જ્યારે
ભાષાની વ્યાખ્યામાં પરિણામનું સૂત્ર $!B$ છે. એ સંજ્ઞા
જાળવવા $[!A]!B$ મૂક્યું છે.}

\begin{figure}
  \begin{center}
    \begin{tikzpicture}[modal]
      \node[world] (w1) [label={below left:$p, \lnot q$}]{$w_1$}; 
      \node[world] (w2) [left=of w1, label={left:$\lnot p, \lnot q$}]{$w_2$}; 
      \node[world] (w3) [above=of w1, label={right:$p, q$}] {$w_3$};
      \draw[<->] (w1) to [swap] node {$b$} (w2);
      \draw[->,loop below] (w1) to node {$a, b$} (w1);
      \draw[<->] (w1) to [swap] node {$a$} (w3);
      \draw[->,loop above] (w2) to node {$a, b$} (w2) ;
      \draw[->,loop above] (w3) to node {$a, b$} (w3) ;
      
      \node[world](v1)[right of=w1, xshift=1.25in, label={right:$p, \lnot q$}]{$w'_1$};
      \node[world](v3)[above of=v1,label={right:$p,q$}]{$w'_3$};
      \draw[<->] (v1) to node {$a$} (v3);
      \draw[->,loop below] (v1) to node {$a,b$} (v1);
      \draw[->,loop above] (v3) to node {$a,b$} (v3) ;
      
      \node(m)[below=of w1]{$\mModel M$};
      \node(m')[below=of v1]{$\mModel M \mid p$}; 

      \draw[-,dotted] (w1) to node {\footnotesize $p$ની જાહેરાત}(v1) ;
     
    \end{tikzpicture}
  \end{center}
  \caption{$p$ની જાહેર જાહેરાત પહેલાં અને પછી.}
  \ollabel{fig:announcement-example}
\end{figure}

!!a{formula}ની જાહેર જાહેરાતને નિદર્શનું સંકોચન, એટલે કે
!!{formula} સત્ય હોય તે જ જગતો સુધી તેને મર્યાદિત કરવું,
એમ સમજી શકીએ છીએ. \olref{fig:announcement-example}માં~$p$ની
જાહેર જાહેરાતની અસરનું ઉદાહરણ છે. તેમાં નોંધવા જેવી વાત એ છે
કે જાહેરાતના પરિણામે કર્તા~$b$ને~$p$ની ખબર પડે છે, જ્યારે
કર્તા~$a$ને નવી ખબર પડતી નથી (કારણ કે $a$ને~$p$ સત્ય હોવાની
પહેલેથી જ ખબર હતી).

વધુ ઔપચારિક રીતે, $\mSat{M}{\lnot \Knows_b p}[w_1]$ છે,
પરંતુ $\mSat{M \mid p}{\Knows_b p}[w'_1]$ છે. તેથી
$\mSat{M}{[p] \Knows_b p}[w_1]$. જાહેરાતના પરિણામ વિશે
એથી પણ પ્રબળ દાવો કરી શકાય છે. હકીકતમાં,
$\mSat{M}{[p]\CKnows_{\{a,b\}} p}[w_1]$ પણ છે. બીજા શબ્દોમાં,
$p$ જાહેર થયા પછી તે \emph{સામાન્ય જ્ઞાન} બને છે.

પણ શું સામાન્ય રીતે પણ એમ જ થાય? શું~$!A$ની સત્ય જાહેરાત
\emph{હંમેશા} $!A$ને સામાન્ય જ્ઞાન બનાવે છે? કદાચ આશ્ચર્યની વાત
છે કે જવાબ ના છે. હકીકતમાં, અમુક !!{formula}s સત્યપૂર્વક
જાહેર કરી શકાય, પરંતુ જાહેરાત પછી તે સત્ય રહેતાં નથી. દાખલા
તરીકે, \olref{fig:announcement-example}માં~$w_1$ ખાતે
$p \land \lnot \Knows_b p$ની જાહેરાતની અસર વિચારો.
\sourcecorrection{OLMD-082}{મૂળમાં $\mModel M \mid p$ અને
$\mModel M \mid (p \land \lnot \Knows_b p)$ એક જ નિદર્શ
હોવાનું કહે છે; આ આકૃતિમાં પહેલામાં બે જગતો અને બીજામાં
એક જ જગત બચે છે. છતાં બંનેમાં મૂળ પ્રથમ જગતમાં જાહેરાત
પછી $b$ને $p$ ખબર પડે છે, તેથી મુખ્ય નિષ્કર્ષ જળવાય છે.}
હકીકતમાં, $\mModel M \mid p$ અને
$\mModel M \mid (p \land \lnot \Knows_b p)$ એક જ નિદર્શ
નથી: બીજામાં ફક્ત પ્રથમ જગત બચે છે, જ્યારે પહેલામાં પ્રથમ
અને ત્રીજું બંને બચે છે. જોકે આપણે અગાઉ નોંધ્યું છે તેમ,
$\mSat{M \mid p}{\Knows_b p}[w'_1]$. તેથી
$\mSat{M \mid (p \land \lnot \Knows_b p)}{\lnot
(p \land \lnot \Knows_b p)}[w'_1]$. આમ આ !!a{formula}
જાહેર થયા પછી અસત્ય બને છે.

\end{document}
